\section{Objet de l'annexe}
\label{annexe-resistance-objet}

Cette annexe développe les calculs qui fondent la triade

\[
\text{inertie}
\longrightarrow
\text{résistance}
\longrightarrow
\text{appropriation}
\]

dans le modèle continu R0.

Elle traite séparément :

\begin{itemize}
\item
  la perte de stabilité du repos ;
\item
  la naissance et le développement de la branche résistante ;
\item
  la stabilité de cette branche ;
\item
  la rupture de stabilité de type Hopf ;
\item
  le caractère non structurel du seuil \(g_{\mathrm{ign}}\) ;
\item
  la définition de l'appropriation par les bassins ;
\item
  la construction du prédicteur de la durée minimale
  \(T^\ast(g)\).
\end{itemize}

Les calculs analytiques seront distingués des résultats obtenus par
continuation ou intégration numérique.

\section{Système forcé réduit}
\label{annexe-resistance-systeme}

Le système source \(A\) est fixé sur son puit positif. Sa projection sur le système receveur \(B\) vaut

\[
w=\tanh(x_A)>0.
\]

Pour les paramètres de référence,

\[
w=0{,}98006.
\]
Cette valeur est la projection mesurée du puits positif du système
source complet utilisé dans les expériences. Elle ne doit pas être
recalculée à partir de la racine arrondie \(x\simeq2{,}30\) du système
réduit autonome, même si les deux valeurs sont numériquement très
proches.

Sur la diagonale interne de \(B\),

\[
h_B=(x,x),
\qquad
v_B=(v,v),
\qquad
\tau_B=(\theta,\theta).
\]

La norme de la trace est

\[
r=\sqrt2\,|\theta|.
\]

La dynamique forcée est

\[
\dot x=v,
\]

\[
\dot v
=
-\Gamma(r)v
-
m(r)x
+
\zeta\tanh x,
\]

\[
\dot\theta
=
\alpha\tanh x
+
gw
-
\beta\theta,
\]

avec

\[
m(r)
=
\frac{1+r^2}
{1+(1+\lambda)r^2}.
\]

Dans tout le voisinage de la branche étudiée,

\[
\theta>0,
\]

donc

\[
r=\sqrt2\,\theta.
\]

\section{Le repos exact reste invariant}
\label{annexe-resistance-repos-invariant}

Supposons

\[
x(0)=0,
\qquad
v(0)=0.
\]

Alors, quelle que soit la valeur de \(\theta(t)\),

\[
\dot x=0,
\]

et

\[
\dot v
=
-\Gamma(\sqrt2|\theta|)\,0
-
m(\sqrt2|\theta|)\,0
+
\zeta\tanh0
=
0.
\]

Par unicité,

\[
x(t)=0,
\qquad
v(t)=0
\]

pour tout \(t\).

La trace suit alors l'équation autonome

\[
\dot\theta
=
gw-\beta\theta.
\]

Sa solution est

\[
\theta(t)
=
e^{-\beta t}\theta(0)
+
\frac{gw}{\beta}
\left(
1-e^{-\beta t}
\right).
\]

Ainsi,

\[
\theta(t)
\longrightarrow
\theta_0(g)
=
\frac{gw}{\beta}.
\]

\begin{quote}
\resultheading{Proposition}{prop:D-repos-invariant}{}
L'ensemble
\[
\{x=0,\;v=0\}
\]
est invariant pour toute intensité \(g\). L'influence charge la trace,
mais elle ne crée pas elle-même un déplacement.
\end{quote}

\section{Perte de stabilité du repos}
\label{annexe-resistance-gc}

Autour de

\[
x=0,
\qquad
v=0,
\qquad
\theta=\theta_0(g),
\]

la dynamique de \((x,v)\) est

\[
\frac{\mathrm d}{\mathrm dt}
\begin{pmatrix}
x\\ v
\end{pmatrix}
=
\begin{pmatrix}
0&1\\
\zeta-m(r_0(g))&-\Gamma(r_0(g))
\end{pmatrix}
\begin{pmatrix}
x\\ v
\end{pmatrix},
\]

avec

\[
r_0(g)
=
\sqrt2\,\theta_0(g)
=
\frac{\sqrt2\,gw}{\beta}.
\]

Le polynôme caractéristique est

\[
s^2
+
\Gamma(r_0(g))s
+
m(r_0(g))-\zeta.
\]

La friction est positive. Le repos est donc localement stable si

\[
m(r_0(g))>\zeta,
\]

et instable si

\[
m(r_0(g))<\zeta.
\]

Le seuil de trace défini dans l'annexe précédente vérifie

\[
m(\tau_{\mathrm{ign}})=\zeta.
\]

Ainsi,

\[
r_0(g_c)=\tau_{\mathrm{ign}},
\]

d'où

\[
\boxed{
g_c
=
\frac{\beta\tau_{\mathrm{ign}}}
{\sqrt2\,w}.
}
\]

Avec

\[
\tau_{\mathrm{ign}}
=
0{,}377964\ldots,
\]

on obtient

\[
\boxed{
g_c
=
0{,}03273\ldots
}
\]

pour les paramètres de référence.

Le troisième mode, porté par la trace, possède toujours la valeur
propre \(-\beta\). Le changement de stabilité du repos vient donc bien
du bloc \((x,v)\).

\section{Équation de la branche d'équilibres}
\label{annexe-resistance-residu}

À l'équilibre,

\[
v^\ast=0,
\]

et

\[
\theta^\ast(x,g)
=
\frac{\alpha\tanh x+gw}{\beta}.
\]

Posons

\[
r(x,g)
=
\sqrt2\,\theta^\ast(x,g)
=
\frac{\sqrt2}{\beta}
\left(
\alpha\tanh x+gw
\right).
\]

L'équation restante est

\[
m(r(x,g))x
-
\zeta\tanh x
=
0.
\]

On définit le résidu

\[
\boxed{
F(x,g)
=
m(r(x,g))x
-
\zeta\tanh x.
}
\]

Comme

\[
F(0,g)=0
\]

pour tout \(g\), l'origine forme une branche d'équilibres.

Pour \(x\neq0\), on peut diviser par \(x\) et définir

\[
G(x,g)
=
m(r(x,g))
-
\zeta\frac{\tanh x}{x}.
\]

La fonction \(\tanh x/x\) est prolongée continûment en \(x=0\) par la
valeur \(1\). On obtient ainsi une fonction lisse au voisinage de
\((0,g_c)\).

Les équilibres non nuls sont les solutions de

\[
G(x,g)=0.
\]

\section{Dérivées utiles}
\label{annexe-resistance-derivees-m}

La fonction

\[
m(r)
=
\frac{1+r^2}
{1+(1+\lambda)r^2}
\]

vérifie

\[
\boxed{
m'(r)
=
-\frac{2\lambda r}
{\left[
1+(1+\lambda)r^2
\right]^2}.
}
\]

Pour \(r>0\),

\[
m'(r)<0.
\]

Sa dérivée seconde est

\[
\boxed{
m''(r)
=
-\frac{
2\lambda
\left[
1-3(1+\lambda)r^2
\right]
}{
\left[
1+(1+\lambda)r^2
\right]^3
}.
}
\]

Nous utiliserons les notations

\[
r_c=\tau_{\mathrm{ign}},
\]

\[
m_1=m'(r_c),
\qquad
m_2=m''(r_c),
\]

\[
A=\frac{\sqrt2\,\alpha}{\beta},
\qquad
B=\frac{\sqrt2\,w}{\beta},
\]

et

\[
\delta=g-g_c.
\]

\section{Naissance transcritique de la résistance}
\label{annexe-resistance-transcritique}

Au voisinage de \((x,\delta)=(0,0)\),

\[
\tanh x
=
x-\frac{x^3}{3}+O(x^5),
\]

donc

\[
\frac{\tanh x}{x}
=
1-\frac{x^2}{3}+O(x^4).
\]

Par ailleurs,

\[
r(x,g)
=
r_c
+
Ax
+
B\delta
+
O(x^3).
\]

Au premier ordre,

\[
\begin{aligned}
G(x,g)
&=
m_1(Ax+B\delta)
+
O\left(
x^2+|x\delta|+\delta^2
\right).
\end{aligned}
\]

Ainsi,

\[
\frac{\partial G}{\partial x}(0,g_c)
=
m_1A
\neq0.
\]

Le théorème des fonctions implicites donne une branche unique

\[
x=x_{\mathrm{res}}(g)
\]

de solutions non nulles de \(G(x,g)=0\) passant par

\[
x_{\mathrm{res}}(g_c)=0.
\]

Sa pente est

\[
x_{\mathrm{res}}'(g_c)
=
-
\frac{G_g(0,g_c)}
{G_x(0,g_c)}
=
-
\frac{m_1B}{m_1A}.
\]

Donc

\[
\boxed{
x_{\mathrm{res}}'(g_c)
=
-\frac{w}{\alpha}.
}
\]

Comme \(w>0\) et \(\alpha>0\),

\[
x_{\mathrm{res}}(g)<0
\]

pour \(g>g_c\) suffisamment proche de \(g_c\).

Le résidu complet s'écrit alors

\[
F(x,g)
=
x\,G(x,g).
\]

La branche \(x=0\) et la branche
\(x=x_{\mathrm{res}}(g)\) se croisent en \(g=g_c\), avec échange de
stabilité. Il s'agit de la forme locale d'une bifurcation
transcritique \cite{kuznetsov2004elements}.

\begin{quote}
\resultheading{Proposition}{prop:D-bifurcation-transcritique}{}
La branche résistante naît en \(g=g_c\) par une bifurcation
transcritique. Pour \(g>g_c\) proche du seuil, elle est
antiparallèle à l'influence :
\[
x_{\mathrm{res}}(g)<0.
\]
\end{quote}

\section{Développement au premier ordre}
\label{annexe-resistance-premier-ordre}

Le calcul précédent donne immédiatement

\[
\boxed{
x_{\mathrm{res}}(g)
=
-\frac{w}{\alpha}
(g-g_c)
+
O\left(
(g-g_c)^2
\right).
}
\]

Pour les paramètres de référence,

\[
-\frac{w}{\alpha}
=
-2{,}17791\ldots
\]

La pente vient d'une compensation simple. Près de \(g_c\),

\[
\theta^\ast
=
\frac{
\alpha x_{\mathrm{res}}
+
gw
}{\beta}
+
O(x_{\mathrm{res}}^3).
\]

En remplaçant

\[
x_{\mathrm{res}}
\approx
-\frac{w}{\alpha}(g-g_c),
\]

on obtient

\[
\alpha x_{\mathrm{res}}+gw
\approx
-w(g-g_c)+gw
=
g_cw.
\]

La contre-écriture interne neutralise donc, au premier ordre, l'excès
de charge externe au-dessus du seuil critique.

\section{Développement au second ordre}
\label{annexe-resistance-second-ordre}

Pour obtenir le terme quadratique, il faut conserver tous les termes
de degré total deux dans \(G\).

On a

\[
G(x,g)
=
a_1\delta
+
a_2x
+
a_{xx}x^2
+
a_{xg}x\delta
+
a_{gg}\delta^2
+
O\left(
(|x|+|\delta|)^3
\right),
\]

avec

\[
a_1=m_1B,
\]

\[
a_2=m_1A,
\]

\[
a_{xx}
=
\frac12m_2A^2
+
\frac{\zeta}{3},
\]

\[
a_{xg}
=
m_2AB,
\]

\[
a_{gg}
=
\frac12m_2B^2.
\]

Cherchons

\[
x_{\mathrm{res}}(g)
=
c_1\delta
+
c_2\delta^2
+
O(\delta^3).
\]

À l'ordre un,

\[
a_1+a_2c_1=0,
\]

d'où

\[
c_1=-\frac{B}{A}=-\frac{w}{\alpha}.
\]

À l'ordre deux,

\[
a_2c_2
+
a_{xx}c_1^2
+
a_{xg}c_1
+
a_{gg}
=
0.
\]

Ainsi,

\[
\boxed{
c_2
=
-
\frac{
a_{xx}c_1^2
+
a_{xg}c_1
+
a_{gg}
}{
a_2
}.
}
\]

Pour les paramètres de référence,

\[
a_{xx}
=
5{,}77917\ldots,
\]

\[
a_{xg}
=
24{,}01147\ldots,
\]

\[
a_{gg}
=
26{,}14742\ldots
\]

Les trois contributions au numérateur valent approximativement

\[
27{,}40,
\qquad
-52{,}31,
\qquad
26{,}15.
\]

Elles s'annulent presque. On obtient

\[
c_2
\simeq
0{,}3220
\]

à la précision retenue dans les calculs numériques.

Finalement,

\[
\boxed{
x_{\mathrm{res}}(g)
=
-2{,}1774(g-g_c)
+
0{,}3220(g-g_c)^2
+
O\left(
(g-g_c)^3
\right).
}
\]

La quasi-linéarité de la branche ne vient donc pas d'un oubli du terme
quadratique. Elle vient d'une compensation entre trois contributions
distinctes.

\section{La trace n'est pas exactement épinglée au seuil}
\label{annexe-resistance-marginalite}

À un équilibre non nul,

\[
m(r^\ast)x^\ast
=
\zeta\tanh x^\ast.
\]

En divisant par \(x^\ast\),

\[
\boxed{
m(r^\ast)
=
\zeta\frac{\tanh x^\ast}{x^\ast}.
}
\]

Pour \(x^\ast\neq0\),

\[
\left|
\frac{\tanh x^\ast}{x^\ast}
\right|
<1.
\]

Ainsi,

\[
m(r^\ast)<\zeta.
\]

Près de l'origine,

\[
\frac{\tanh x}{x}
=
1-\frac{x^2}{3}+O(x^4),
\]

donc

\[
\boxed{
m(r^\ast)
=
\zeta
-
\frac{\zeta}{3}(x^\ast)^2
+
O\left(
(x^\ast)^4
\right).
}
\]

La résistance ne maintient pas exactement le système sur la marge
\(m=\zeta\). Elle s'en écarte d'un terme quadratique.

\section{Jacobien sur une branche forcée}
\label{annexe-resistance-jacobien}

Soit un équilibre

\[
(x^\ast,0,\theta^\ast),
\qquad
\theta^\ast>0.
\]

Posons

\[
r^\ast=\sqrt2\,\theta^\ast,
\]

\[
\Gamma_\ast=\Gamma(r^\ast),
\]

\[
s_\ast=\operatorname{sech}^2x^\ast,
\]

\[
c_\ast=\zeta s_\ast,
\]

\[
\kappa_\ast
=
m(r^\ast)-c_\ast,
\]

et

\[
b_\ast
=
-\sqrt2\,m'(r^\ast)x^\ast.
\]

Le Jacobien est

\[
\boxed{
J_\ast
=
\begin{pmatrix}
0&1&0\\
-\kappa_\ast&-\Gamma_\ast&b_\ast\\
\alpha s_\ast&0&-\beta
\end{pmatrix}.
}
\]

Son polynôme caractéristique est

\[
p(s)
=
s^3+a_2s^2+a_1s+a_0,
\]

avec

\[
\boxed{
a_2
=
\Gamma_\ast+\beta,
}
\]

\[
\boxed{
a_1
=
\Gamma_\ast\beta+\kappa_\ast,
}
\]

\[
\boxed{
a_0
=
\beta\kappa_\ast
-
\alpha b_\ast s_\ast.
}
\]

Les conditions de Routh--Hurwitz pour un polynôme cubique sont

\[
a_2>0,
\qquad
a_1>0,
\qquad
a_0>0,
\qquad
a_2a_1>a_0.
\]

\section{Signe des coefficients sur la branche résistante}
\label{annexe-resistance-signe-coefficients}

Sur la branche résistante,

\[
x^\ast<0,
\]

et

\[
m'(r^\ast)<0.
\]

Ainsi,

\[
b_\ast
=
-\sqrt2\,m'(r^\ast)x^\ast
<0.
\]

Par ailleurs, l'équation d'équilibre donne

\[
m(r^\ast)
=
\zeta\frac{\tanh x^\ast}{x^\ast}.
\]

Donc

\[
\kappa_\ast
=
\zeta
\left(
\frac{\tanh x^\ast}{x^\ast}
-
\operatorname{sech}^2x^\ast
\right).
\]

Pour \(x\neq0\),

\[
\frac{\tanh x}{x}
>
\operatorname{sech}^2x.
\]

En effet, pour \(x>0\), cette inégalité équivaut à

\[
\tanh x
>
x\,\operatorname{sech}^2x,
\]

déjà démontrée lors de l'étude de la fonction
\(x/\tanh x\). Les deux membres sont pairs après division par \(x\),
donc le résultat vaut aussi pour \(x<0\).

Ainsi,

\[
\kappa_\ast>0.
\]

On obtient alors

\[
a_2>0,
\qquad
a_1>0,
\]

et

\[
a_0
=
\beta\kappa_\ast
+
\alpha|b_\ast|s_\ast
>0.
\]

La seule condition susceptible de changer de signe est donc

\[
a_2a_1-a_0.
\]

Un calcul direct donne

\[
\boxed{
a_2a_1-a_0
=
\Gamma_\ast^2\beta
+
\Gamma_\ast\beta^2
+
\Gamma_\ast\kappa_\ast
+
\alpha b_\ast s_\ast.
}
\]

Comme \(b_\ast<0\), le dernier terme est négatif. Il traduit le fait
que, sur une branche antiparallèle, la boucle
trace--état peut déstabiliser l'équilibre par oscillation.

Près de \(g_c\), le terme négatif est d'ordre \(|x^\ast|\), tandis que
\(\kappa_\ast\) est d'ordre \((x^\ast)^2\). Les termes positifs
\(\Gamma_\ast^2\beta+\Gamma_\ast\beta^2\) restent non nuls. La branche
est donc stable immédiatement après sa naissance.

\section{Condition de Hopf}
\label{annexe-resistance-hopf-condition}

Supposons

\[
a_2>0,
\qquad
a_1>0,
\qquad
a_0>0.
\]

Lorsque

\[
a_2a_1=a_0,
\]

le polynôme se factorise :

\[
p(s)
=
(s+a_2)(s^2+a_1).
\]

Les valeurs propres sont alors

\[
-a_2,
\qquad
\pm i\sqrt{a_1}.
\]

La pulsation de Hopf est donc

\[
\boxed{
\omega_H=\sqrt{a_1}.
}
\]

Pour conclure à une bifurcation de Hopf, il faut encore vérifier que la
paire complexe traverse l'axe imaginaire avec une vitesse non nulle.
Cette condition de transversalité est évaluée numériquement le long de
la branche \cite{kuznetsov2004elements}.

\section{Localisation numérique de la rupture}
\label{annexe-resistance-hopf-numerique}

La branche d'équilibres est poursuivie numériquement en \(g\). Pour
chaque point, on calcule les trois valeurs propres du Jacobien.

Le changement de signe est observé entre

\[
g=0{,}120
\]

et

\[
g=0{,}132.
\]

La localisation affinée donne

\[
\boxed{
g_H
=
0{,}1285.
}
\]

À ce point,

\[
|x_{\mathrm{res}}(g_H)|
\simeq
0{,}209,
\]

et

\[
\boxed{
\omega_H
\simeq
0{,}334,
\qquad
\frac{2\pi}{\omega_H}
\simeq
18{,}8.
}
\]

La partie réelle dominante vaut approximativement

\[
-0{,}0025
\]

à \(g=0{,}120\),

\[
0
\]

à \(g=0{,}1285\), et

\[
+0{,}0009
\]

à \(g=0{,}132\).

Le statut de ces valeurs est \texttt{[PROP-NUM]} :

\begin{itemize}
\item
  la condition de Hopf est analytique ;
\item
  la valeur \(g_H\) est obtenue par continuation numérique ;
\item
  la transversalité est vérifiée numériquement ;
\item
  le premier coefficient de Lyapunov n'est pas encore calculé.
\end{itemize}

La branche ne disparaît pas en \(g_H\). Par exemple, à

\[
g=0{,}145,
\]

elle possède encore une racine

\[
x_{\mathrm{res}}
=
-0{,}2453,
\]

mais sa valeur propre dominante a une partie réelle positive.

Il est donc préférable de parler de \emph{rupture de stabilité} plutôt
que de disparition de la résistance.

\section{Trois seuils de nature différente}
\label{annexe-resistance-trois-seuils}

Les valeurs

\[
g_c,
\qquad
g_{\mathrm{ign}},
\qquad
g_H
\]

ne décrivent pas le même phénomène.

\subsection{\texorpdfstring{$g_c$}{gc} : bifurcation locale}

Le seuil

\[
g_c
=
0{,}0327
\]

est défini par une égalité spectrale :

\[
m(r_0(g_c))=\zeta.
\]

Il marque l'échange de stabilité entre le repos et la branche
résistante.

\subsection{\texorpdfstring{$g_H$}{gH} : rupture de stabilité de la branche}

Le seuil

\[
g_H
=
0{,}1285
\]

est défini par la condition

\[
a_2a_1=a_0.
\]

Il marque la rupture locale de stabilité de la branche résistante par franchissement d'une paire complexe.

\subsection{\texorpdfstring{$g_{\mathrm{ign}}$}{gign} : frontière de bassin}

Fixons une condition initiale et un protocole. Pour la condition
standard,

\[
x_0=10^{-3},
\qquad
v_0=0,
\qquad
\theta_0=0,
\]

et une mise sous contact brutale, on observe

\[
\boxed{
g_{\mathrm{ign}}
\simeq
0{,}0556.
}
\]

Cette valeur ne correspond à aucune perte d'existence ou de stabilité
de la branche résistante.

À

\[
g=g_{\mathrm{ign}}+\varepsilon,
\]

une condition initiale placée exactement sur la branche résistante y
reste, tandis que la condition initiale standard s'en échappe.

Le seuil doit donc être écrit

\[
\boxed{
g_{\mathrm{ign}}
=
g_{\mathrm{ign}}
(x_0,v_0,\theta_0;\mathcal P),
}
\]

où \(\mathcal P\) désigne le protocole de mise sous contact.

La hiérarchie observée est

\[
\boxed{
g_c
<
g_{\mathrm{ign}}
(x_0,v_0,\theta_0;\mathcal P)
<
g_H.
}
\]

\section{Influence du protocole}
\label{annexe-resistance-protocole}

La dépendance à \(\mathcal P\) peut être montrée sans changer la cible
finale \(g\).

À

\[
g=0{,}10,
\]

trois protocoles ont été comparés :

\begin{itemize}
\item
  mise sous contact brutale ;
\item
  rampe courte de durée \(20\) ;
\item
  rampe lente de durée \(400\).
\end{itemize}

Les deux premiers conduisent à l'appropriation. La rampe lente permet
à la trajectoire de suivre la branche résistante stable et conduit à
la résistance.

Ce résultat ne constitue pas un théorème adiabatique. Il montre
numériquement que \(g_{\mathrm{ign}}\) est une frontière de capture,
sensible à l'histoire du chargement.

\section{Définition de l'appropriation par les bassins}
\label{annexe-resistance-definition-appropriation}

Soit

\[
\Phi_{g,\mathcal P}^T
\]

l'application qui associe à l'état initial \(z_0\) l'état obtenu après
une interaction d'intensité \(g\), de durée \(T\), selon le protocole
\(\mathcal P\).

Après la fin de l'interaction, on revient à la dynamique autonome. Pour
un état \(z\), notons

\[
\omega_0(z)
\]

son ensemble limite sous cette dynamique.

\begin{quote}
\resultheading{Définition}{def:D-appropriation}{}
L'interaction
\[
(g,T,\mathcal P)
\]
est appropriée depuis \(z_0\) si
\[
\omega_0
\left(
\Phi_{g,\mathcal P}^T(z_0)
\right)
=
\mathcal A_1,
\]
\[
\omega_0(z_0)
=
\mathcal A_0,
\]
avec
\[
\mathcal A_1\neq\mathcal A_0,
\]
où \(\mathcal A_0\) et \(\mathcal A_1\) sont des attracteurs autonomes
du système receveur.
\end{quote}

Cette définition ne demande pas que les attracteurs soient des points
fixes. Elle s'applique aussi aux régimes hybrides invariants décrits
dans l'annexe précédente.

Elle exclut :

\begin{itemize}
\item
  une réponse transitoire qui disparaît après découplage ;
\item
  une copie maintenue uniquement par la présence du système source ;
\item
  une transmission directe du support de mémoire du système source.
\end{itemize}

\section{Durée minimale d'exposition}
\label{annexe-resistance-duree-minimale}

Pour une intensité \(g\), une condition initiale \(z_0\) et un
protocole \(\mathcal P\), on définit

\[
\boxed{
T^\ast
(g;z_0,\mathcal P)
=
\inf
\left\{
T>0:
(g,T,\mathcal P)
\text{ est appropriée depuis }z_0
\right\}.
}
\]

Lorsque l'ensemble est vide,

\[
T^\ast=+\infty.
\]

La notation abrégée

\[
T^\ast(g)
\]

suppose donc que la condition initiale et le protocole sont fixés.

\section{Trajectoire de charge utilisée dans le prédicteur}
\label{annexe-resistance-charge-lineaire}

Le prédicteur est construit près du repos. Dans l'équation de la
trace, on néglige d'abord la contribution interne

\[
\alpha\tanh x.
\]

Pour une exposition rectangulaire de durée \(T\), avec

\[
\theta(0)=0,
\]

la charge approchée est

\[
\theta_{g,T}^{(0)}(t)
=
\frac{gw}{\beta}
\left(
1-e^{-\beta t}
\right)
\]

pour

\[
0\leq t\leq T.
\]

Après découplage,

\[
\theta_{g,T}^{(0)}(t)
=
\frac{gw}{\beta}
\left(
1-e^{-\beta T}
\right)
e^{-\beta(t-T)}
\]

pour

\[
t\geq T.
\]

La norme correspondante est

\[
r_{g,T}^{(0)}(t)
=
\sqrt2
\left|
\theta_{g,T}^{(0)}(t)
\right|.
\]

\section{Système linéaire non autonome}
\label{annexe-resistance-systeme-lineaire}

On linéarise la dynamique de \((x,v)\) le long de cette charge :

\[
\frac{\mathrm d}{\mathrm dt}
\begin{pmatrix}
x_{\mathrm{lin}}\\
v_{\mathrm{lin}}
\end{pmatrix}
=
A_{g,T}(t)
\begin{pmatrix}
x_{\mathrm{lin}}\\
v_{\mathrm{lin}}
\end{pmatrix},
\]

avec

\[
\boxed{
A_{g,T}(t)
=
\begin{pmatrix}
0&1\\
\zeta-m(r_{g,T}^{(0)}(t))
&
-\Gamma(r_{g,T}^{(0)}(t))
\end{pmatrix}.
}
\]

La condition initiale est

\[
x_{\mathrm{lin}}(0)=x_0,
\qquad
v_{\mathrm{lin}}(0)=v_0.
\]

Notons

\[
x_u>0
\]

la position de l'équilibre-selle autonome positif sur la diagonale.
Pour les paramètres de référence,

\[
x_u=0{,}0717.
\]

Le prédicteur considère que l'appropriation devient possible lorsque
la trajectoire linéarisée atteint ce seuil avant que la charge ne soit
oubliée.

Une définition opérationnelle sans temps de coupure arbitraire est

\[
\boxed{
T_{\mathrm{lin}}^\ast(g)
=
\inf
\left\{
T>0:
\sup_{t\geq0}
x_{\mathrm{lin}}(t;g,T)
\geq
x_u
\right\}.
}
\]

Le supremum est fini dans les cas étudiés, car la charge finit par
décroître et le repos autonome redevient stable.

\section{Statut du prédicteur}
\label{annexe-resistance-statut-predicteur}

La construction précédente n'est pas une preuve d'équivalence entre le
système linéaire et la dynamique non linéaire complète.

Elle repose sur trois approximations :

\begin{enumerate}
\item
  la contribution interne \(\alpha\tanh x\) est négligée dans la
  charge de référence ;
\item
  la dynamique de \((x,v)\) est linéarisée autour du repos ;
\item
  le franchissement de \(x_u\) est utilisé comme approximation du
  franchissement de la séparatrice complète.
\end{enumerate}

Le résultat a donc le statut \texttt{[PROP-NUM]}.

Il possède néanmoins deux propriétés importantes :

\begin{itemize}
\item
  aucun paramètre n'est ajusté sur les valeurs de \(T^\ast\) ;
\item
  le seuil \(x_u\) vient du paysage autonome, calculé indépendamment.
\end{itemize}

\section{Comparaison avec la dynamique complète}
\label{annexe-resistance-comparaison-Tstar}

Dans la suite, le protocole de référence est l'exposition rectangulaire à
mise sous contact brutale.

Pour la condition initiale standard et le protocole rectangulaire, on
obtient :

\[
\begin{array}{c|cc}
g
&
T_{\mathrm{lin}}^\ast(g)
&
T^\ast(g)\text{ mesuré}
\\
\hline
0{,}05
&
+\infty
&
+\infty
\\
0{,}08
&
31{,}6
&
31{,}7
\\
0{,}10
&
24{,}8
&
24{,}8
\\
0{,}15
&
17{,}8
&
17{,}9
\\
0{,}20
&
14{,}9
&
15{,}0
\\
0{,}30
&
12{,}2
&
12{,}3
\end{array}
\]

L'écart relatif reste inférieur à \(0{,}5\,\%\) pour les valeurs
finies testées.

Dans la zone où la trajectoire est capturée par la branche résistante,
le prédicteur classe correctement

\[
T^\ast=+\infty.
\]

Cette concordance montre que la frontière d'appropriation est
principalement gouvernée par trois ingrédients :

\[
\boxed{
\text{charge de la trace}
+
\text{instabilité temporaire du repos}
+
\text{franchissement de la selle autonome}.
}
\]

\section{Pourquoi prolonger l'exposition ne suffit pas toujours}
\label{annexe-resistance-prolonger-exposition}

Pour

\[
g_c<g<g_{\mathrm{ign}}
\]

selon le protocole de référence, la branche résistante est stable.

Une exposition plus longue ne rapproche donc pas nécessairement la
trajectoire du bassin du puits positif. Elle laisse au contraire davantage de
temps au transitoire pour être capturé par l'équilibre résistant.

Après découplage, la charge externe disparaît. La branche résistante
forcée se déplace puis cesse d'exister comme équilibre du système
autonome. La trajectoire retourne alors vers le repos.

On obtient ainsi une situation où

\[
\text{durée de contact}\to\infty
\]

n'implique pas

\[
\text{appropriation}.
\]

Cette propriété distingue la résistance d'un simple retard
d'allumage.

\section{Réductions exactes sous couplage mutuel}
\label{annexe-resistance-couplage-mutuel}

Considérons deux systèmes identiques, couplés symétriquement.

\subsection{Sous-espace aligné}

Si

\[
h_B=h_A=h,
\]

alors

\[
\tanh(h_B)=\tanh(h_A).
\]

L'équation de trace devient

\[
\dot\tau
=
(\alpha+g)\tanh h
-
\beta\tau.
\]

Le couplage est donc exactement équivalent au remplacement

\[
\boxed{
\alpha_{\mathrm{eff}}
=
\alpha+g.
}
\]

Cette réduction est valable seulement dans le sous-espace aligné
invariant et pour des systèmes identiques avec couplage symétrique.

\subsection{Sous-espace antipodal}

Si

\[
h_B=-h_A=h,
\]

alors

\[
\tanh(h_B)=-\tanh(h_A).
\]

L'équation de trace devient

\[
\dot\tau
=
(\alpha-g)\tanh h
-
\beta\tau.
\]

On obtient

\[
\boxed{
\alpha_{\mathrm{eff}}
=
\alpha-g.
}
\]

Lorsque \(g>\alpha\), le coefficient effectif devient négatif. Le
sous-espace reste formellement invariant, mais l'interprétation comme
écriture interne renforcée n'est plus applicable de la même manière.

\section{Seuil de compétition}
\label{annexe-resistance-seuil-competition}

Pour une valeur

\[
k_{\mathrm{eff}}
=
\frac{\alpha_{\mathrm{eff}}}{\beta},
\]

l'équation d'un équilibre actif à \(q\) composantes non nulles est

\[
x
=
\zeta
\Lambda_q(x;k_{\mathrm{eff}})
\tanh x,
\]

avec

\[
\Lambda_q(x;k)
=
1+
\frac{
\lambda qk^2\tanh^2x
}{
1+qk^2\tanh^2x
}.
\]

Pour \(x>0\), cette équation équivaut à

\[
\zeta
=
\frac{x}
{
\Lambda_q(x;k_{\mathrm{eff}})
\tanh x
}.
\]

Définissons

\[
\boxed{
\zeta_{\min}(k)
=
\min_{x>0}
\frac{x}
{
\Lambda_q(x;k)\tanh x
}.
}
\]

Un état actif existe si

\[
\zeta\geq\zeta_{\min}(k).
\]

Dans le sous-espace antipodal,

\[
k
=
\frac{\alpha-g}{\beta}.
\]

Le seuil de compétition vérifie donc

\[
\boxed{
\zeta_{\min}
\left(
\frac{\alpha-g_{\mathrm{comp}}}{\beta}
\right)
=
\zeta.
}
\]

Pour les paramètres de référence et la branche à deux composantes
actives, la résolution numérique donne

\[
\boxed{
g_{\mathrm{comp}}
=
0{,}3757.
}
\]

Sous ce seuil, les états antipodaux existent encore. Au-dessus, ils ne
sont plus disponibles.

La formule est analytique ; la valeur numérique vient de la
minimisation et de la résolution d'une équation scalaire.

\section{Résumé des statuts}
\label{annexe-resistance-resume-statuts}

Les résultats se répartissent ainsi.

\subsection{Résultats analytiques}

\begin{itemize}
\item
  invariance du repos exact ;
\item
  formule de \(g_c\) ;
\item
  existence locale de la branche résistante ;
\item
  caractère transcritique de sa naissance ;
\item
  pente
  \[
  -w/\alpha;
  \]
\item
  formule du terme quadratique ;
\item
  correction de la marginalité exacte ;
\item
  Jacobien et conditions de Routh--Hurwitz ;
\item
  condition algébrique de Hopf ;
\item
  définition de l'appropriation par les bassins ;
\item
  réductions exactes
  \[
  \alpha_{\mathrm{eff}}=\alpha\pm g.
  \]
\end{itemize}

\subsection{Résultats analytiques complétés numériquement}

\begin{itemize}
\item
  coefficient quadratique
  \[
  0{,}3220;
  \]
\item
  localisation
  \[
  g_H=0{,}1285;
  \]
\item
  pulsation
  \[
  \omega_H\simeq0{,}334;
  \]
\item
  seuil
  \[
  g_{\mathrm{comp}}=0{,}3757;
  \]
\item
  prédicteur de
  \[
  T^\ast(g).
  \]
\end{itemize}

\subsection{Faits de bassin ou de protocole}

\begin{itemize}
\item
  valeur
  \[
  g_{\mathrm{ign}}\simeq0{,}0556
  \]
  pour la condition initiale standard ;
\item
  protection par une rampe lente ;
\item
  éjection vers le puits positif après la rupture de stabilité.
\end{itemize}

\section{Limites}
\label{annexe-resistance-limites}

Cette annexe ne démontre pas :

\begin{itemize}
\item
  la sous-criticité ou la sur-criticité de la bifurcation de Hopf ;
\item
  la géométrie globale de la séparatrice ;
\item
  une formule fermée pour
  \(g_{\mathrm{ign}}\) ;
\item
  l'exactitude du prédicteur linéaire pour toutes les conditions
  initiales ;
\item
  la persistance de la triade complète hors de la diagonale interne ;
\item
  la stabilité des branches sous couplage asymétrique.
\end{itemize}

Le résultat nécessaire à la suite est néanmoins établi :

\[
\boxed{
\begin{minipage}{0.86\linewidth}
L'influence continue ne crée pas directement un mouvement. Elle
déplace la stabilité du repos. Au premier seuil, une branche
antiparallèle stable apparaît et produit une résistance. Une condition
initiale peut quitter son bassin avant la rupture structurelle de cette
branche : c'est le seuil de capture \(g_{\mathrm{ign}}\). Lorsque la
branche perd finalement sa stabilité par Hopf, elle persiste comme
configuration mais ne peut plus retenir la trajectoire.
\end{minipage}
}
\]
